\section{可视证据与研究时间线}

\subsection{Slide 3：Q-Learning 模型管线}

这一页展示 Q-Learning from data 到 deployment 的 pipeline，每个环节都需要 guardrail（replay hygiene、target network、threshold monitor）。

\begin{figure}[H]
\centering
\includegraphics[width=0.92\textwidth,page=3]{06_cs224r_qlearning_2025.pdf}
\caption{Slide：Q-Learning pipeline and guardrail checks}
\end{figure}
\newpage

\subsection{Slide 4：Temporal difference 与 replay dynamics}

Slide 4 把 TD error 与 replay sampling 关联起来，强调 `high-error` 采样对 prioritized buffer 的贡献同时也要人工审查。

\begin{figure}[H]
\centering
\includegraphics[width=0.92\textwidth,page=4]{06_cs224r_qlearning_2025.pdf}
\caption{Slide：TD error monitoring 与 prioritized sampling}
\end{figure}
\newpage

\subsection{Slide 5：Deployment 监控}

部署时需要对 reward distribution、episode success rate、latency 进行多通道监控，Slide 5 展示了 telemetry stack。

\begin{figure}[H]
\centering
\includegraphics[width=0.92\textwidth,page=5]{06_cs224r_qlearning_2025.pdf}
\caption{Slide：deployment telemetry and alert stack}
\end{figure}
\newpage

\subsection{Slide 6：Action timeline and evidence matrix}

Slide 6 附上 evidence matrix，说明每种 failure 都配有对应 action timeline，正如我们在本文中添加的 table 所示。

\begin{figure}[H]
\centering
\includegraphics[width=0.92\textwidth,page=6]{06_cs224r_qlearning_2025.pdf}
\caption{Slide：Evidence matrix and action timeline}
\end{figure}
\newpage

\subsection{Slide 7：Research directions}

这页点出 future direction：multi-step Q-Learning, safety evaluation, sim-to-real.

\begin{figure}[H]
\centering
\includegraphics[width=0.92\textwidth,page=7]{06_cs224r_qlearning_2025.pdf}
\caption{Slide：Research directions for safe Q-Learning}
\end{figure}
\newpage

\subsection{Slide 8：Time/Loss visualization}

Slide 8 展示了 reward curve 与 TD error 的同步下降，为收束全文提供视觉线索。

\begin{figure}[H]
\centering
\includegraphics[width=0.92\textwidth,page=8]{06_cs224r_qlearning_2025.pdf}
\caption{Slide：Reward & TD error visualization}
\end{figure}
\newpage

